\noindent\textbf{Implications for Deep Research.} These findings motivate the Deep Research (DR) follow-up explored in RQ3 (Section~\ref{sec:rq3_results}). The key insight is that high retrieval coverage does not imply strong causal evidence: in complex biomedical domains (e.g., Uleman's), standard web scraping retrieves many sources but the returned evidence is often correlational, conditional, or hedged with respect to the specific edge claim, yielding low citation-judge scores. By contrast, in physics-based CLDs with less causal ambiguity, when citations are retrievable they more often contain direct causal statements and receive high citation-judge scores. Together with reasonable human-validation alignment for the citation judge, this suggests weak discrimination in biomedical settings is driven primarily by evidence type and domain/evidence dependence (what evidence exists and is practically retrievable), rather than a failure of judge functioning.

However, before we can conclude that, one alternative explanation remains: the quality of retrieved evidence. In the validation set, we fed the judge known-to-be-correct references to isolate its functioning, but it may be that in our current search process evidence of this quality is not being found due to the limits of the naive search our citation judge employs. The Citation Judge (through the citation fetcher and search API components) receives topically relevant literature by querying the causal narrative verbatim (up to an arbitrary 5 sources), but this is a naive approach and does not specifically target direct causal evidence (RCTs, experimental manipulations, meta-analyses) or use alternative query formulations using field-specific terminology. Deep Research addresses this limitation through: (1) multiple parallel search queries targeting causal mechanisms rather than topical keywords; (2) iterative evidence gathering using refined queries based on retrieved evidence when initial results are inconclusive; (3) a multi-agent pipeline that scores, selects, and judges evidence for causal specificity; and (4) explicit filtering for study types that establish causation rather than correlation.

